pdfoutput=1
\pdfoutput=1
% !TEX program = xelatex
\documentclass[12pt,a4paper]{article}

% ================================================================
% Packages
% ================================================================
\usepackage{amsmath,amssymb,amsthm}
\usepackage{mathtools}
\usepackage{bm}
\usepackage{graphicx}
\usepackage{hyperref}
\usepackage{geometry}
\usepackage{enumitem}
\usepackage{booktabs}
\usepackage{tikz}
\usepackage{tikz-cd}
\usepackage{xcolor}
\usepackage{cleveref}
\usepackage{bbm}
\usepackage{float}
\usepackage{algorithm}
\usepackage{algpseudocode}

\geometry{margin=2.5cm}

% ================================================================
% Hyperref setup
% ================================================================
\hypersetup{
    colorlinks=true,
    citecolor=blue,
    urlcolor=blue,
    pdftitle={Potential Surface Misalignment in Medical Diagnosis: Gauge Misalignment Between the Doctor's and Patient's Coordinate Systems},
    pdfauthor={SCX},
    pdfsubject={SCX Medical Gauge Theory},
}

% ================================================================
% Theorem environments
% ================================================================
\newtheorem{theorem}{Theorem}[section]
\newtheorem{lemma}[theorem]{Lemma}
\newtheorem{proposition}[theorem]{Proposition}
\newtheorem{corollary}[theorem]{Corollary}
\newtheorem{definition}[theorem]{Definition}
\newtheorem{conjecture}[theorem]{Conjecture}
\newtheorem{assumption_env}[theorem]{Assumption}
\newtheorem{criterion}[theorem]{Criterion}

\theoremstyle{remark}
\newtheorem{remark}[theorem]{Remark}
\newtheorem{example}[theorem]{Example}
\newtheorem{protocol}[theorem]{Protocol}

% ================================================================
% Custom commands
% ================================================================
\newcommand{\R}{\mathbb{R}}
\newcommand{\C}{\mathbb{C}}
\newcommand{\N}{\mathbb{N}}
\newcommand{\E}{\mathbb{E}}
\newcommand{\Pbb}{\mathbb{P}}
\newcommand{\Var}{\mathrm{Var}}
\newcommand{\Cov}{\mathrm{Cov}}
\newcommand{\loss}{\mathcal{L}}
\newcommand{\D}{\mathcal{D}}
\newcommand{\G}{\mathcal{G}}
\newcommand{\T}{\mathcal{T}}
\newcommand{\SCX}{\textsf{SCX}}
\newcommand{\med}{\textsf{med}}
\newcommand{\pt}{\textsf{pt}}
\newcommand{\gaugeGroup}{\mathbb{G}}

\DeclareMathOperator{\argmax}{arg\,max}
\DeclareMathOperator{\argmin}{arg\,min}
\DeclareMathOperator{\sgn}{sgn}
\DeclareMathOperator{\diag}{diag}
\DeclareMathOperator{\spec}{spec}
\DeclareMathOperator{\softmax}{softmax}
\DeclareMathOperator{\Diagnose}{Diagnose}
\DeclareMathOperator{\Misdiagnose}{Misdiagnose}

% Honest critique
\newcommand{\honestcrit}[1]{%
    \textcolor{red}{\textbf{[Honest Strike]}#1}%
}

% Proof-status markers
\newcommand{\rigorFull}{\textnormal{\textbf{[Rigorous Proof]}}}
\newcommand{\rigorPartial}{\textnormal{\textbf{[Partial Proof]}}}
\newcommand{\rigorSketch}{\textnormal{\textbf{[Proof Sketch]}}}
\newcommand{\openproblem}{\textnormal{\textbf{[Open Problem]}}}

% Assumption tags
\newcommand{\assumptionRef}[1]{\textbf{A#1}}

% ================================================================
% Title
% ================================================================
\title{\textbf{Potential Surface Misalignment in Medical Diagnosis}\\[4pt]
\large Gauge Misalignment Between the Doctor's and Patient's Coordinate Systems}
\author{SCX}
\date{July 1, 2026}

\begin{document}

\maketitle

\begin{center}
\footnotesize
\textbf{Version:} v1.0 \quad | \quad
\textbf{Status:} Preprint \quad | \quad
\textbf{Classification:} SCX Theoretical System --- Medical Volume: Diagnostic Standardization
\end{center}

% ================================================================
% Abstract
% ================================================================

\begin{abstract}

The core operation of medical diagnosis---a physician listens to a patient's symptom description and produces a diagnosis---is implicitly assumed to be a direct comparison process: the physician matches the patient's described clinical presentation against known disease patterns. This paper proves that this default assumption is wrong. Physicians and patients live in \textbf{different coordinate systems}: patients describe symptoms in their subjective experiential coordinate system (quality of pain, location of discomfort, degree of fatigue), while physicians interpret symptoms in their medical training coordinate system (pathological classification, laboratory reference ranges, imaging features). No natural, objective translation relationship exists between the two---this is a \textbf{gauge freedom}. When one coordinate system is transformed into another, the comparison itself is mathematically ill-defined---this is precisely the \textbf{deep mathematical root cause of misdiagnosis}.

We introduce the SCX potential surface geometric framework to formalize this problem. The patient's subjective state defines a potential surface $\mathcal{S}_{\pt}$ (energy distribution in symptom space), and the physician's knowledge system defines another potential surface $\mathcal{S}_{\med}$ (energy distribution in disease space). When the two coordinate systems are misaligned---i.e., $\mathbf{g}_{\med} \neq \mathbf{g}_{\pt}$---the physician's interpretation of the patient's symptoms is a \textbf{gauge-dependent} operation whose conclusion depends on undeclared coordinate system choices. We prove: (1) single-physician diagnosis, under the framework of Theorem~3 (Honest Person Theorem), cannot distinguish genuine symptoms from patient narrative noise---this is a hard boundary on diagnostic uncertainty; (2) multi-expert consultation ($M > 1$) provides a \textbf{gauge-fixing} mechanism: the consensus detection of $M$ independent observers reduces misdiagnosis at the exponential convergence rate $e^{-2M\Delta^2}$ (the medical translation of Theorem~1); (3) patient autonomy is not an ethical appeal---it is a mathematical necessity: $\sum \mathbf{g}_m = \mathbf{0}$ demands that physicians must not default their medical coordinate system as the sole standard, and the patient's coordinate system must participate in the gauge-fixing process; (4) the optimal role of AI-assisted diagnosis is not to replace physicians---it is to serve as a \textbf{potential surface auditor}, detecting misalignment between physician and patient coordinate systems without holding its own coordinate system.

We develop a complete operational framework: (a) a pre-consultation coordinate system declaration protocol---physician and patient each declare their reference frame; (b) gauge-fixing through multi-expert consultation---using the independent coordinate systems of $M$ physicians to construct a common reference plane; (c) attitude auditing---detecting whether a physician harbors undeclared attitudinal bias through repeated interactions (corollary: attitude leaks through repeated interactions); (d) AI potential surface auditing---using neural networks to compute the gradient consistency of the physician's diagnostic trajectory in the patient's symptom space, identifying gauge misalignment. The paper contains 7 theorems/propositions, 3 operational protocols, and 1 Honest Strike section.

\vspace{4pt}
\noindent\textbf{Keywords:} medical diagnosis, gauge theory, potential surface misalignment, doctor-patient communication, multi-expert consultation, diagnostic error, patient autonomy, AI audit, SCX framework

\end{abstract}

% ================================================================
\section*{Core Intuition: Understanding Coordinate System Misalignment in One Page}
\addcontentsline{toc}{section}{Core Intuition: Understanding Coordinate System Misalignment in One Page}
% ================================================================

\vspace{-8pt}

\begin{center}
\fbox{%
\begin{minipage}{0.92\textwidth}

\subsection*{The Interpreter Metaphor}

Imagine a French patient seeing a Chinese physician. The two share no common language.

The patient describes their suffering in French: \textit{``J'ai une douleur lancinante dans le bas du dos, \c{c}a irradie vers la jambe gauche, et c'est pire le matin.''}

The physician must diagnose. But the physician does not understand French. Two paths lie ahead:

\textbf{Path A (status quo):} The physician hires an interpreter. The interpreter translates French into Chinese. But the interpreter has their own understanding---some French words have no precise Chinese equivalent (\textit{lancinante} = throbbing? stabbing? pulling?). The interpreter chooses one translation. The physician makes a diagnosis based on this translation. If the translation is wrong---the diagnosis is wrong.

\textbf{Path B (multiple observers):} The physician hires three interpreters. Each independently translates the same French passage. If all three produce the same translation---the physician can trust it. If the three produce three different versions---the physician knows they face an unreliable signal and should not base a diagnosis on any single translation.

\textbf{This is potential surface misalignment in diagnosis.}

\begin{itemize}
    \item \textbf{French = the patient's coordinate system} $\mathcal{S}_{\pt}$. The patient feels pain in their own body---this experience is real, complete, and requires no translation.
    \item \textbf{Chinese = the physician's coordinate system} $\mathcal{S}_{\med}$. The physician understands disease within their knowledge system---pathology, imaging, laboratory indicators.
    \item \textbf{Translation = the gauge transformation} $\mathbf{g}_{\med \to \pt}$. After hearing the patient's description, the physician finds a correspondence in their own coordinate system. This correspondence is not unique---different physicians (different interpreters) will make different correspondences.
    \item \textbf{Translation error = misdiagnosis.} Not because the physician is unskilled---but because no predefined gauge-fixing exists between the two coordinate systems.
    \item \textbf{Three interpreters = multi-expert consultation} $M=3$. Theorem~1: the consensus of $M$ independent observers converges to the truth at exponential rate.
\end{itemize}

\textbf{The solution is not hiring more interpreters---it is aligning coordinate systems first.} Let the patient write it down in French (coordinate system declaration). Let the physician list their understood French-Chinese correspondence table (gauge-fixing). Then---align before comparing. Or more cleverly---do not translate. Let the patient evaluate their own status changes in their coordinate system, let the physician evaluate objective indicators in their coordinate system, and let both coordinate systems operate independently, validating consensus only in \textbf{output space} (outcome metrics: whether pain decreased, whether function recovered, whether biochemical markers improved). Output space is the only space that requires no gauge-fixing---because reduced pain is real in the patient's coordinate system, and improved biochemical markers are real in the physician's coordinate system.

\end{minipage}%
}
\end{center}

\vspace{8pt}

\begin{center}
\fbox{%
\begin{minipage}{0.92\textwidth}

\subsection*{Three-Sentence Summary}

\begin{enumerate}
    \item \textbf{Problem identified:} Misdiagnosis is not (only) about insufficient knowledge or negligence---it is about physicians and patients living in two different coordinate systems, with no predefined gauge-fixing between them. The physician is comparing the incomparable.
    \item \textbf{Solution approaches:} Three pathways. (a) Pre-consultation coordinate system declaration---physician and patient each clarify their reference frame. (b) Multi-expert consultation---the consensus of $M>1$ independent physicians eliminates misdiagnosis at the exponential rate $e^{-2M\Delta^2}$. (c) Validation in output space---do not compare intermediate processes, only compare final outcomes, because outcomes are real in each coordinate system independently.
    \item \textbf{Surprise finding:} Patient autonomy is not an ethical concept---it is a mathematical necessity. $\sum\mathbf{g}_m = \mathbf{0}$ requires the patient's coordinate system to be incorporated into the gauge-fixing process. Healthcare systems that fail to respect patient autonomy mathematically produce higher misdiagnosis rates.
\end{enumerate}

\end{minipage}%
}
\end{center}

\vspace{12pt}

% ================================================================
\section{Introduction: From Potential Surface Geometry to Medical Diagnosis}
\label{sec:intro}
% ================================================================

In the SCX Theoretical System, potential surface geometry is the unified mathematical language for understanding multi-observer epistemology. Each observer defines a potential surface $\mathcal{S}$---a function mapping ``what is right,'' ``what is important,'' and ``what deserves attention'' onto energy/probability. Differences between the potential surfaces of multiple observers arise naturally---not because anyone is ``wrong,'' but because training data, prior beliefs, and observational perspective cause each to independently converge to different local minima.

In prior work, we have identified and addressed potential surface misalignment in the following domains:
\begin{itemize}
    \item \textbf{Atomic potential function merging:} ACE potential functions for different chemical elements exhibit gauge freedom in shared correction terms---identical physical predictions can arise from different parameter configurations. Gauge-fixing via post-processing orthogonal projection reduces gauge violation to machine precision [EGP paper].
    \item \textbf{MoE routing:} Different experts in multi-expert networks define gauge-inequivalent representations in their respective output spaces---the router compares the incomparable. MILP gauge-fixing restores routing optimality, and SVD spectrum detection identifies shared hallucinations [MoE gauge paper].
    \item \textbf{Data quality auditing:} Label noise and sample intrinsic difficulty are indistinguishable under a single observer (Honest Person Theorem). Multi-expert consensus detects noise at exponential convergence rate (Theorem~1), and this rate is precisely the constant that minimizes the maximum [SCX theory paper].
\end{itemize}

\honestcrit{Medicine is the next target. And the most important target---because here potential surface misalignment does not (only) affect model accuracy, it affects human lives.}

\subsection{The Medical Motivation of This Paper: Why Misdiagnosis Is Not ``Individual Physician Error''}

Globally, deaths and severe harm caused by diagnostic error each year are significant [WHO Global Patient Safety Report]. Traditional medical education attributes misdiagnosis to:
\begin{enumerate}[label=(\arabic*)]
    \item Insufficient physician knowledge---never learned about this rare disease
    \item Physician negligence---failed to ask the key question, failed to order the key test
    \item Unclear patient communication---patient unable to describe symptoms accurately
    \item Time pressure---outpatient visit too short for adequate history-taking
\end{enumerate}

These explanations share a common assumption: \textbf{there exists a ``correct'' diagnosis that can be obtained from the patient's symptoms by a ``sufficiently good'' physician.} Misdiagnosis is a deviation from this ideal---the physician is not good enough, the patient is not clear enough, the time is not long enough.

\textbf{The core claim of this paper is: this assumption is wrong.}

It is not that ``the physician is not good enough, hence misdiagnosis''---it is that \textbf{even a perfect physician under perfect conditions faces an ineliminable mathematical boundary on single-physician diagnosis.} This boundary does not come from the limits of knowledge---it comes from the \textbf{incomparability of coordinate systems}.

\subsection{Core Concept: Misalignment of Two Potential Surfaces}

\begin{definition}[The Patient's Symptom Potential Surface]\label{def:patient_surface}
Let the patient's subjective symptom space be $\Omega_{\pt} \subset \R^{d_{\pt}}$, where each dimension represents a symptom dimension perceivable by the patient (pain location, pain quality, pain intensity, accompanying symptoms, temporal pattern, etc.). The patient's state defines a potential function:
\begin{equation}
    \mathcal{S}_{\pt}: \Omega_{\pt} \to \R, \quad \mathcal{S}_{\pt}(\omega) = \text{the patient's subjective distress level in state }\omega
\end{equation}
The patient maps their current potential state $\mathcal{S}_{\pt}(\omega_0)$ to a symptom description $s \in \Sigma$ ($\Sigma$ is the linguistic description space) through language.
\end{definition}

\begin{definition}[The Physician's Disease Potential Surface]\label{def:doctor_surface}
Let the physician's medical knowledge space be $\Omega_{\med} \subset \R^{d_{\med}}$, where each dimension represents a diagnostic dimension in the physician's cognition (pathological classification, laboratory indicators, imaging features, epidemiological probability, treatment response patterns, etc.). The physician's knowledge system defines a potential function:
\begin{equation}
    \mathcal{S}_{\med}: \Omega_{\med} \to \R, \quad \mathcal{S}_{\med}(\omega) = \text{the ``match''/``likelihood'' in the physician's cognition of state }\omega\text{ corresponding to a disease}
\end{equation}
The physician maps the patient's description $s$ to a diagnosis $d \in \mathcal{D}$ ($\mathcal{D}$ is the disease classification space) through the diagnostic operation.
\end{definition}

\textbf{Core problem:} These two potential surfaces are defined on different spaces ($\Omega_{\pt} \neq \Omega_{\med}$), with different dimensions ($d_{\pt} \neq d_{\med}$), and no predefined mapping exists between them.

\begin{definition}[Diagnostic Gauge Transformation]\label{def:diagnostic_gauge}
The diagnostic operation $F: \Sigma \to \mathcal{D}$ consists of two steps:
\begin{enumerate}[label=(\roman*)]
    \item \textbf{Listening-interpretation} (gauge transformation): $T_{\mathbf{g}}: \Sigma \to \Omega_{\med}$, where $\mathbf{g} \in \G$ is the physician's gauge parameter. The physician maps the patient's linguistic description $s$ to a point $T_{\mathbf{g}}(s) \in \Omega_{\med}$ in the physician's medical space.
    \item \textbf{Pattern matching} (diagnosis): $\mathcal{S}_{\med}(T_{\mathbf{g}}(s))$---the physician evaluates the energy/likelihood of the mapped position on their potential surface.
\end{enumerate}
The complete diagnostic function is:
\begin{equation}
    d = F_{\mathbf{g}}(s) = \argmin_{d \in \mathcal{D}} \left\| \mathcal{S}_{\med}(T_{\mathbf{g}}(s)) - \mathcal{S}_{\med}(\omega_d) \right\|
\end{equation}
where $\omega_d$ is the ``prototype'' position of disease $d$ in $\Omega_{\med}$.
\end{definition}

\honestcrit{The gauge transformation $T_{\mathbf{g}}$ is the invisible core of diagnosis. The physician does not consciously ``choose'' a gauge---the gauge is implicitly determined by their training history, clinical experience, cultural background, and even their mood that day. Two equally excellent physicians, facing the same patient, produce two different diagnoses---not because one is wrong, but because their gauges $\mathbf{g}_1$ and $\mathbf{g}_2$ differ.}

\subsection{Formal Definition of Misdiagnosis}

\begin{definition}[Misdiagnosis as Potential Surface Misalignment]\label{def:misdiagnosis}
Let the patient's true state be $\omega_{\pt}^* \in \Omega_{\pt}$, with corresponding true disease (if present) $d^* \in \mathcal{D}$. The physician's diagnosis is $d = F_{\mathbf{g}}(s)$, where $s$ is the patient's linguistic description of $\omega_{\pt}^*$. Then:
\begin{equation}
    \Misdiagnose(d, d^*) = \mathbf{1}\{d \neq d^*\}
\end{equation}

The root cause of misdiagnosis is \textbf{potential surface misalignment}:
\begin{equation}
    \Pbb(\Misdiagnose) \propto \|T_{\mathbf{g}}(s) - \omega_{\pt \to \med}^*\|
\end{equation}
where $\omega_{\pt \to \med}^*$ is the ``correct'' projection of the patient's true state into the physician's space---but this projection \textbf{does not exist} unless the two coordinate systems have already been aligned.
\end{definition}

\textbf{Ill-defined!} The definition of $\omega_{\pt \to \med}^*$ itself depends on a gauge---whose gauge? The patient's? The physician's? A hypothetical ``God's perspective''? This paper proves: in the absence of explicit gauge-fixing, this ``correct projection'' mathematically does not exist---it is a gauge-dependent quantity whose value differs under different gauges.

% ================================================================
\section{Formalization of Diagnostic Gauge Freedom}
\label{sec:gauge_formal}
% ================================================================

\subsection{The Gauge Group in Medical Diagnosis}

\begin{definition}[Medical Diagnosis Gauge Group]\label{def:medical_gauge_group}
Let the implicit gauge parameter used by physician $m$ in the diagnostic operation be $\mathbf{g}_m \in \G$. The gauge group $\G$ contains the following substructures:

\begin{enumerate}[label=(\roman*)]
    \item \textbf{Translation gauge} $\G^{\text{trans}} = \{\mathbf{b} \in \R^{d_{\med}}\}$: shifts the patient's linguistic description to different positions in the medical space. $\mathbf{b} = 0$ means ``I believe the literal meaning of the patient's description''; $\mathbf{b} \neq 0$ means ``I think the patient's description needs adjustment---perhaps exaggeration, perhaps understatement, perhaps cultural expression differences.''
    
    \item \textbf{Scaling gauge} $\G^{\text{scale}} = \{\alpha > 0\}$: amplifies or diminishes the severity of the patient's symptoms. $\alpha = 1$ is literal correspondence; $\alpha > 1$ is ``I think this is more severe than the patient says''; $\alpha < 1$ is ``I think the patient is exaggerating.''
    
    \item \textbf{Prior gauge} $\G^{\text{prior}} = \{\bm{\pi} \in \Delta^{|\mathcal{D}|-1}\}$: the physician's prior belief about disease prevalence. This prior is not explicit---it comes from the physician's training history, practice region, and specialty orientation.
    
    \item \textbf{Risk gauge} $\G^{\text{risk}} = \{r \in [0,1]\}$: the physician's degree of risk aversion. $r \to 1$ means ``better to over-diagnose than to miss'' (defensive medicine); $r \to 0$ means ``better to watchfully wait than to over-diagnose'' (minimal intervention principle).
\end{enumerate}

The full diagnostic gauge group is the product of these subgroups:
\begin{equation}
    \G = \G^{\text{trans}} \times \G^{\text{scale}} \times \G^{\text{prior}} \times \G^{\text{risk}}
\end{equation}
\end{definition}

\honestcrit{Note: there is no ``zero gauge.'' Claiming ``I listened objectively to the patient's description'' is itself a gauge choice---$\mathbf{b}=0, \alpha=1$---but the legitimacy of this choice is unverified. The patient may indeed have exaggerated ($\alpha$ should be $<1$), or may indeed have underreported ($\alpha$ should be $>1$). ``Objectivity'' is not the absence of a gauge---it is an undeclared gauge.}

\begin{theorem}[Gauge Non-Invariance of Single-Physician Diagnosis]\label{thm:noninvariant}\rigorFull
Let two physicians $D_1$ and $D_2$ face the same patient description $s$. Suppose the two physicians possess identical medical knowledge $\mathcal{S}_{\med}$ (perfect training), but hold different gauge parameters $\mathbf{g}_1 \neq \mathbf{g}_2$. Then there exists a patient description $s$ such that $F_{\mathbf{g}_1}(s) \neq F_{\mathbf{g}_2}(s)$---i.e., two perfect physicians can produce different diagnoses.
\end{theorem}

\begin{proof}
We demonstrate how gauge differences produce diagnostic disagreement. Let the patient description $s$ lie in the ``border zone'' between two disease prototypes $\omega_{d_A}$ and $\omega_{d_B}$ in $\Sigma$. Under gauge $T_{\mathbf{g}_1}$, $s$ is mapped to a point in $\Omega_{\med}$ closer to $\omega_{d_A}$; under gauge $T_{\mathbf{g}_2}$, $s$ is mapped to another point closer to $\omega_{d_B}$. Whenever both $\mathbf{g}_1 \neq \mathbf{g}_2$ and $s$ lies in the boundary region between two diseases hold simultaneously, diagnostic disagreement necessarily occurs.

More rigorously: define the diagnostic boundary
\begin{equation}
    \mathcal{B}(d_A, d_B) = \{\omega \in \Omega_{\med} : \mathcal{S}_{\med}(\omega - \omega_{d_A}) = \mathcal{S}_{\med}(\omega - \omega_{d_B})\}
\end{equation}
The gauge transformation $T_{\mathbf{g}}$ maps points from the patient description space $\Sigma$ into $\Omega_{\med}$. If there exists $s$ such that $T_{\mathbf{g}_1}(s)$ and $T_{\mathbf{g}_2}(s)$ fall on opposite sides of the diagnostic boundary $\mathcal{B}(d_A, d_B)$, then the diagnoses differ. Since $T_{\mathbf{g}}$ is continuous and $\mathbf{g}_1 \neq \mathbf{g}_2$, such $s$ exists in a region of nonzero measure.
\end{proof}

\begin{corollary}[Diagnostic Disagreement Is Not ``Error''---It Is Gauge Difference]\label{cor:disagreement}
When two qualified physicians produce different diagnoses for the same patient, this does not necessarily mean at least one physician ``made a mistake.'' It may mean the two physicians used different gauge parameters to process the same patient description. Before accusing either physician of ``misdiagnosis,'' one must first fix the gauge---i.e., first determine what the ``correct gauge'' is.
\end{corollary}

\honestcrit{But who defines the ``correct gauge''? The physician community? The patient? Clinical guidelines? This question has no answer within medicine---it is a coordinate system selection problem, and coordinate system selection is a political-ethical operation, not a medical operation.}

\subsection{The Medical Consultation as an Insufficiently Gauge-Declared Operation}

\begin{definition}[Standard Consultation Process---Insufficient Gauge Declaration]\label{def:standard_visit}
The standard consultation process can be formalized as:
\begin{enumerate}[label=(\arabic*)]
    \item Patient enters the consultation room---patient coordinate system $\mathcal{S}_{\pt}$ implicit but undeclared
    \item Physician takes the history---physician maps the patient's answers into $\Omega_{\med}$ under the physician's own gauge $\mathbf{g}_{\med}$
    \item Physical examination---physician samples within $\Omega_{\med}$ (auscultation, palpation, inspection)
    \item Ancillary testing---physician obtains high-precision measurements in $\Omega_{\med}$ via laboratory/imaging
    \item Diagnosis---physician searches $\mathcal{S}_{\med}$ for the lowest-energy disease match
\end{enumerate}

\textbf{Note step 2:} The patient's coordinate system $\mathcal{S}_{\pt}$ and the physician's gauge $\mathbf{g}_{\med}$ are \textbf{never explicitly aligned} in this workflow. The physician measures the patient's drawing with their own ruler---but the zero mark on the ruler and the origin on the drawing are not the same point.
\end{definition}

\begin{protocol}[Coordinate System Declaration Protocol---Pre-Consultation Gauge-Fixing]\label{prot:coordinate_declaration}
Before diagnosis begins, both parties execute the following declaration operations:
\begin{enumerate}[label=(\arabic*),itemsep=3pt]
    \item \textbf{Patient coordinate system declaration:} The patient declares their reference frame in symptom space---how pain is quantified (0--10 scale or ``no pain/some pain/very severe pain''), temporal patterns of symptoms, self-definition of ``normal.''
    \item \textbf{Doctor coordinate system declaration:} The physician declares their diagnostic framework---the differential diagnosis list for this consultation, the prior probabilities of considered diseases, the interpretation criteria for tests.
    \item \textbf{Gauge-fixing negotiation:} Both parties confirm the intersection points of their coordinate system declarations---in which dimensions the patient can accept the physician's quantification framework, and in which dimensions the physician adopts the patient's subjective description.
    \item \textbf{Difference logging:} Explicitly record the portions where the two coordinate systems cannot be aligned---e.g., the patient insists ``my pain is 'tearing,' not 'stabbing','' but the physician's diagnostic framework only has ``stabbing'' and ``throbbing'' as options. This difference itself is a diagnostic signal.
\end{enumerate}
\end{protocol}

% ================================================================
\section{The Honest Person Theorem---Hard Boundary of Single-Physician Diagnosis}
\label{sec:honest_person_medicine}
% ================================================================

\subsection{Medical Translation of Theorem 3}

The \textbf{Honest Person Theorem} (Theorem~3) from the SCX Theoretical System has an extremely important medical translation:

\begin{theorem}[Medical Honest Person Theorem---Single Physician Cannot Distinguish Genuine Symptoms from Narrative Noise]\label{thm:honest_medical}
For any patient description $s$ and any single physician $D$'s diagnostic operation $F_{\mathbf{g}}$, there exist two data-generating processes (World A and World B) such that:
\begin{enumerate}[label=(\roman*)]
    \item \textbf{World A (genuine organic disease):} The patient truly has disease $d^*$, and symptoms $s$ are a genuine expression of disease $d^*$. The physician, due to gauge misalignment or knowledge limitations, fails to diagnose $d^*$.
    \item \textbf{World B (functional/psychological symptoms---narrative noise):} The patient has no organic disease---symptoms $s$ are ``symptom noise'' generated by psychological factors, cultural expression, or attentional bias. The physician correctly gives no organic diagnosis.
    \item But in these two worlds, \textbf{all observables of the consultation process (patient description $s$, physician diagnosis $d$, test results $\{t_k\}$) are indistinguishable in joint distribution}.
\end{enumerate}
\end{theorem}

\begin{proof}[Proof Sketch]\rigorSketch
Construct an explicit symmetry between the two worlds. World A: disease $d^*$ produces symptom signal $s_{\text{signal}}$, plus the patient's intrinsic description bias $\varepsilon_{\pt}$ yields $s = s_{\text{signal}} + \varepsilon_{\pt}$. World B: no disease, but the patient's description bias $\varepsilon_{\pt}$ itself happens to have the same distribution as $s_{\text{signal}}$---i.e., $s = \varepsilon_{\pt}$, and $\varepsilon_{\pt} \sim \text{distribution identical to } s_{\text{signal}}$.

A single physician, lacking truth-verification means independent of the description (surgical findings, pathological biopsy, long-term follow-up confirmation), cannot distinguish these two worlds. All information available to the physician passes through the ``channel'' of the patient's description---and the transfer function of this channel is gauge-dependent.
\end{proof}

\begin{remark}
This is the mathematical formulation of the fundamental question ``does the patient actually have a disease or not.'' It explains why, in the absence of independent verification means, single-physician diagnosis has an ineliminable uncertainty---not a knowledge problem, but an information channel problem. The physician can only infer from the patient's description, and the patient's description itself contains an indecomposable mixture of signal and noise.
\end{remark}

\honestcrit{This theorem has a serious Honest Strike against the current healthcare system: a large number of ``functional diagnoses'' and ``psychological diagnoses'' may simply be ``we cannot align this patient's coordinate system, so we categorize it as psychological.'' This is not malice---it is the necessary consequence of Theorem~3. When the physician cannot map the patient's description to any known disease in their own coordinate system, the only remaining explanation is ``the patient's description is not accurate enough'' or ``the description itself is a phenomenon requiring treatment''---but this is mathematically equivalent to ``we have abandoned the effort of gauge alignment.''}

\subsection{How Multi-Expert Consultation Breaks Indistinguishability}

SCX Theorem~1 provides the mechanism to break indistinguishability:

\begin{theorem}[Multi-Expert Consultation Reduces Misdiagnosis---Medical Translation of Theorem~1]\label{thm:multiexpert_medicine}
Let $M$ independent physicians each independently interview and diagnose the same patient. Each physician $m$ holds gauge $\mathbf{g}_m$. Define the multi-expert consistency score:
\begin{equation}
    C(s) = \frac{1}{M} \sum_{m=1}^{M} \mathbf{1}\{F_{\mathbf{g}_m}(s) = d\}
\end{equation}

Under the following conditions:
\begin{enumerate}[label=(\assumptionRef{\arabic*}),itemsep=2pt]
    \item \textbf{Independent interview:} Physicians do not communicate with each other; they independently collect history, independently review images, independently form diagnoses
    \item \textbf{Bounded gauge diversity:} $\|\mathbf{g}_m - \mathbf{g}_{m'}\| \leq G_{\max}$ for all $m, m'$---i.e., physicians' gauges are sufficiently diverse
    \item \textbf{Noise symmetry:} The noise in the patient's description (unclear expression, exaggeration, understatement) affects all physicians symmetrically
\end{enumerate}

Then the consistency score satisfies:
\begin{equation}
    \Pbb(\text{majority diagnosis correct}) \geq 1 - \exp\bigl(-2M \Delta^2\bigr)
\end{equation}
where $\Delta$ is the ``correct diagnosis signal strength''---the proportion of physicians for whom the true disease signal exceeds the diagnostic threshold.

\textbf{Translation:} The more physicians consulted, the closer the probability that the majority diagnosis is correct approaches 1 exponentially.
\end{theorem}

\begin{proof}[Proof Sketch]\rigorSketch
Apply Hoeffding's inequality. Each physician's diagnostic decision can be viewed as a Bernoulli trial---correctly diagnosing with probability $p_m$. When $M$ physicians decide independently, the correctness of majority voting is guaranteed by the Chernoff bound. The key lies in assumption (A1) of independent interviewing---if physicians discuss with each other (destroying independence), the convergence rate degrades from exponential to polynomial or worse.
\end{proof}

\begin{remark}
\textbf{Key insight:} Theorem~1 applies in medicine precisely because different physicians' gauges $\mathbf{g}_m$ differ from one another. Gauge difference is not a bug---it is a feature. If all physicians shared the exact same gauge ($\mathbf{g}_1 = \mathbf{g}_2 = \cdots = \mathbf{g}_M$), then no matter how large $M$ is, diagnosis cannot improve---because all physicians make the same systematic error. Gauge diversity is the source of diagnostic robustness.
\end{remark}

\begin{corollary}[The Mathematical Value of Second Opinions]\label{cor:second_opinion}
A second opinion is not redundant---it is a sample of gauge difference. A physician holding a different gauge may see signals that the first physician missed due to gauge reasons. Define the L2 distance of gauge difference:
\begin{equation}
    D_{12} = \|\mathbf{g}_1 - \mathbf{g}_2\|_2
\end{equation}
The marginal diagnostic information of the second opinion is $I(d_2 ; d^* \mid d_1) \propto D_{12}$---the larger the gauge difference, the greater the information gain of the second opinion.
\end{corollary}

% ================================================================
\section{Patient Autonomy---The Mathematical Necessity of Coordinate System Alignment}
\label{sec:patient_autonomy}
% ================================================================

\subsection{The Meaning of $\sum \mathbf{g}_m = \mathbf{0}$ in the Doctor-Patient Relationship}

The core constraint of the SCX Equality Principle is:
\begin{equation}
    \sum_{m} \mathbf{g}_m = \mathbf{0}
\end{equation}

In the doctor-patient relationship, this translates into an extremely specific operational requirement:

\begin{definition}[Doctor-Patient Gauge Equality Principle]\label{def:equality_principle_medical}
In the diagnostic operation, the physician's medical coordinate system $\mathcal{S}_{\med}$ and the patient's subjective experiential coordinate system $\mathcal{S}_{\pt}$ have \textbf{equal initial weight}---neither party may default their coordinate system as the standard before gauge-fixing. Specifically:
\begin{equation}
    \mathbf{g}_{\med} + \mathbf{g}_{\pt} = \mathbf{0}
\end{equation}
where $\mathbf{g}_{\med}$ is the offset by which the physician projects their coordinate system onto the common reference, and $\mathbf{g}_{\pt}$ is the offset by which the patient projects their coordinate system onto the same reference.
\end{definition}

\textbf{This is not an ethical appeal---it is a mathematical necessity.} If the physician unilaterally declares $\mathbf{g}_{\med} = \mathbf{0}$ (``my coordinate system is the standard''), then $\mathbf{g}_{\pt} = \mathbf{0}$ is also forced, meaning the gauge-fixing is completed entirely by the physician---i.e., the patient's coordinate system is forcibly mapped into the physician's framework. The legitimacy of this mapping is untested---its error is a lower bound on the misdiagnosis rate.

\begin{proposition}[Mathematical Formulation of Patient Autonomy]\label{prop:autonomy}
Let $\mathbf{g}_{\text{fixed}}$ be the fixed gauge (common reference). Patient autonomy is equivalent to the patient's right to participate in the gauge-fixing process:
\begin{equation}
    \Pbb(\mathbf{g}_{\pt} \text{ is incorporated into gauge-fixing}) = 1
\end{equation}
If the patient is excluded from the gauge-fixing process (i.e., the physician unilaterally sets $\mathbf{g}_{\text{fixed}} = \mathbf{g}_{\med}$), then the systematic diagnostic bias $\mathbb{E}[\|F_{\mathbf{g}}(s) - d^*\|] \geq \delta_0 > 0$, where $\delta_0$ is the average Bhattacharyya distance between the two coordinate systems.
\end{proposition}

\begin{proof}
When $\mathbf{g}_{\text{fixed}} = \mathbf{g}_{\med}$, the gauge transformation $T_{\mathbf{g}_{\med}}: \Sigma \to \Omega_{\med}$ is defined entirely by the physician's framework. The patient's subjective experience may lie anywhere in $\Sigma$, but the coordinate axes of $\Omega_{\med}$ are laid out according to the physician's classification system. If the patient's true state has no precise correspondence in the physician's coordinate system (rare disease, atypical presentation, cross-cultural symptom expression), then $T_{\mathbf{g}_{\med}}(s)$ necessarily deviates from the patient's true state---this deviation is systematic and does not vanish as the physician's skill improves.
\end{proof}

\subsection{A Potential Surface Critique of ``Doctor Knows Best''}

\honestcrit{Current medical culture harbors an implicit norm: ``doctor knows best.'' From the perspective of potential surface geometry, this assertion is equivalent to declaring $\mathcal{S}_{\med}$ as the only valid coordinate system---$\mathcal{S}_{\pt}$ is demoted to subordinate, to-be-translated, unreliable. This is a gauge choice---not a truth.}

\begin{proposition}[The High Misdiagnosis Cost of ``Doctor Knows Best'']\label{prop:doctor_knows_best}
Under the ``doctor knows best'' gauge ($\mathbf{g}_{\text{fixed}} = \mathbf{g}_{\med}$), the diagnostic error contains two ineliminable components:
\begin{equation}
    \text{Misdiagnosis rate} \geq \underbrace{\varepsilon_{\text{knowledge}}}_{\text{knowledge limitation}} + \underbrace{\varepsilon_{\text{gauge}}}_{\text{gauge misalignment}}
\end{equation}
where $\varepsilon_{\text{gauge}}$ attains its maximum when $\mathbf{g}_{\text{fixed}} = \mathbf{g}_{\med}$, because the patient's coordinate system is entirely excluded.
\end{proposition}

\subsection{Operational Implementation: The Patient's Role in Gauge-Fixing}

\begin{protocol}[Operational Workflow for Patient Participation in Gauge-Fixing]\label{prot:patient_gauge_fixing}
\begin{enumerate}[label=(\arabic*),itemsep=3pt]
    \item \textbf{Before consultation:} The patient completes a ``My Coordinate System'' questionnaire---describing their pain rating habits (conservative or exaggerating), self-calibration of symptom severity, and moments in prior medical encounters when they felt ``heard'' or ``not heard.''
    \item \textbf{During consultation:} Before forming a diagnosis, the physician first restates to the patient their understanding of the patient's coordinate system---``I understand that by 'stabbing pain' you mean... is that correct?''---this is one iteration of gauge-fixing.
    \item \textbf{At diagnosis delivery:} The physician provides not only the diagnosis, but also its ``translation'' in both coordinate systems:
    \begin{itemize}
        \item In the doctor coordinate system: pathological diagnosis + objective indicators + treatment plan
        \item In the patient coordinate system: symptom explanation (``the X you feel is because of Y'') + expected changes (``after treatment you will feel Z'') + self-monitorable indicators for the patient
    \end{itemize}
    \item \textbf{During follow-up:} Gauge-fixing iterates continuously---the patient provides feedback on subjective changes after treatment, and the physician adjusts gauge parameters based on this feedback.
\end{enumerate}
\end{protocol}

% ================================================================
\section{Attitude Auditing---The Physician's Implicit Gauge Leakage}
\label{sec:attitude_audit}
% ================================================================

The proposition ``attitude leaks through repeated interactions'' (from the SCX Theoretical System) has a direct application in the doctor-patient relationship:

\begin{proposition}[Attitude Auditing---Detectability of Physician Attitude]\label{prop:attitude}
A physician's implicit gauge bias $\mathbf{g}_{\med} - \mathbf{g}_{\text{claimed}}$ (the difference between actual gauge and claimed gauge) can be detected through a series of repeated interactions. Specifically, let a physician, across $T$ independent consultations (facing different patients but similar symptom description patterns), produce a diagnostic sequence $\{d_t\}_{t=1}^{T}$. Then:
\begin{equation}
    \lim_{T \to \infty} \frac{1}{T} \sum_{t=1}^{T} \|d_t - d_t^*\| = \delta_{\mathbf{g}}
\end{equation}
where $\delta_{\mathbf{g}}$ is the systematic diagnostic shift caused by the physician's true gauge bias.

\textbf{Translation:} A physician's attitude (``this patient exaggerates symptoms,'' ``old people just like to complain,'' ``young women often have psychological problems'') is not a private psychological state---it is statistically detectable through the pattern of repeated diagnoses.
\end{proposition}

\begin{proof}[Proof Sketch]\rigorSketch
Consider the physician's diagnostic function $F_{\mathbf{g}_{\text{true}}}(s_t)$, where $\mathbf{g}_{\text{true}} = \mathbf{g}_{\text{claimed}} + \bm{\varepsilon}$. $\bm{\varepsilon}$ is the physician's implicit bias. For ``similar but different'' patients---i.e., $\{s_t\}$ are i.i.d.---the expected diagnostic value $\mathbb{E}[d_t]$ will systematically deviate from the expected value of an unbiased diagnosis $\mathbb{E}[d_t^*]$, with the deviation proportional to $\bm{\varepsilon}$. The law of large numbers guarantees that the sample average converges to this systematic bias.
\end{proof}

\subsection{Dimensions of Attitude Bias}

A physician's implicit gauge bias manifests along the following dimensions:

\begin{enumerate}[label=(\alph*),itemsep=2pt]
    \item \textbf{Gender bias:} The probability that a female patient's myocardial infarction symptoms are attributed to anxiety is systematically higher than for male patients [literature]
    \item \textbf{Age bias:} The probability that an elderly patient's pain is attributed to ``normal aging'' is systematically higher than for younger patients [literature]
    \item \textbf{Racial bias:} Patients of different races use the same pain scales with different habits, but physicians' interpretations are not calibrated for this [literature]
    \item \textbf{Socioeconomic bias:} The probability that a low-income patient's symptoms are attributed to ``psychosocial factors'' is systematically higher than for high-income patients [literature]
    \item \textbf{Language bias:} Non-native speakers use vocabulary to describe symptoms that does not match the physician's expected ``standard description,'' causing information loss
\end{enumerate}

Each type of bias is an instance of $\mathbf{g} \neq \mathbf{0}$---the physician, in their own coordinate system, assigns different ``credibility weights'' (non-uniform distribution of the $\alpha$ parameter) to different types of patients.

\honestcrit{This is not a ``bad doctor'' problem---it is a structural necessity of single-coordinate-system diagnosis. When the diagnostic operation allows undeclared gauge freedom, the physician's implicit beliefs necessarily fill this freedom. The only question is whether these implicit beliefs harm certain patients. They do.}

\subsection{Operational Workflow for Attitude Auditing}

\begin{protocol}[Attitude Auditing---Detecting the Physician's Implicit Gauge Bias]\label{prot:attitude_audit}
\begin{enumerate}[label=(\arabic*),itemsep=3pt]
    \item \textbf{Data collection:} Collect complete records of physician $D$ across $T$ independent consultations (patient demographics, chief complaint, physician diagnosis, final confirmed diagnosis---if available).
    \item \textbf{Control group construction:} Group patients by demographic characteristics (gender, age, race, SES), ensuring that the baseline disease prevalence in each group is known or estimable.
    \item \textbf{Bias detection:} Compute the diagnosis rate difference across groups:
    \begin{equation}
        \Delta_{AB} = \frac{\#\{\text{diagnosed X in group A}\}}{\#\{\text{have X in group A}\}} - \frac{\#\{\text{diagnosed X in group B}\}}{\#\{\text{have X in group B}\}}
    \end{equation}
    If $\Delta_{AB}$ is significantly nonzero and cannot be explained by disease prevalence differences, a gauge bias exists.
    \item \textbf{Feedback and correction:} Feed the detected bias back to the physician---not as ``blame,'' but as ``your gauge has a $\Delta$ offset between group A and group B that you may not be aware of.''
    \item \textbf{Gauge re-fixing:} Once aware of the bias, the physician can consciously adjust gauge parameters---using the same $\alpha$ (scaling gauge) and $\mathbf{b}$ (translation gauge) for groups A and B.
\end{enumerate}
\end{protocol}

% ================================================================
\section{AI-Assisted Diagnosis---The Potential Surface Auditor}
\label{sec:ai_audit}
% ================================================================

\subsection{The Correct Role of AI in Diagnosis}

Current AI medical diagnosis research focuses on ``letting AI produce diagnoses''---i.e., letting AI imitate physicians. This is a directional error. The optimal role of AI is not another physician (yet another coordinate system holder), but a \textbf{potential surface auditor}---a meta-observer that holds no coordinate system of its own and specializes in detecting misalignment between other coordinate systems.

\begin{definition}[AI Potential Surface Auditor]\label{def:ai_auditor}
An AI potential surface auditor $\mathcal{A}$ is a function:
\begin{equation}
    \mathcal{A}: (\Sigma, \mathcal{D}, \mathcal{S}_{\med}, \mathcal{S}_{\pt}) \to \R^+
\end{equation}
It receives the patient description, physician diagnosis, the physician's potential surface description, and the patient's potential surface description, and outputs a \textbf{misalignment score} $\psi \in [0,1]$---measuring the degree of gauge mismatch between the two coordinate systems.

The AI auditor \textbf{does not hold its own coordinate system}---it does not define what the ``correct'' diagnosis is. It only computes whether, given both parties' coordinate systems, the physician's diagnostic trajectory has internal consistency.
\end{definition}

\subsection{Audit Mechanism: Gradient Consistency Detection}

\begin{proposition}[Gradient Consistency as Gauge Alignment Metric]\label{prop:gradient_consistency}
Suppose the AI auditor can compute the ``direction of change'' of the patient's state in the patient's space---i.e., if the patient makes a small change in description (``pain changes from 7 to 8''), how should the physician's diagnosis change. If the physician's diagnostic transformation between two descriptions is inconsistent with the gradient of $\mathcal{S}_{\med}$ in that region, then gauge misalignment or knowledge deficiency exists.

Formally: add a small perturbation $\delta s$ to the patient description $s$ (e.g., change pain from ``throbbing'' to ``stabbing''), the physician's diagnostic change is:
\begin{equation}
    \nabla_s F_{\mathbf{g}}(s) \approx \frac{F_{\mathbf{g}}(s + \delta s) - F_{\mathbf{g}}(s)}{\|\delta s\|}
\end{equation}
If $\nabla_s F_{\mathbf{g}}(s)$ is not smooth in the patient's $\Omega_{\pt}$---i.e., adjacent symptom descriptions produce jump discontinuities in diagnosis---then $\mathbf{g}$ is unfixed in that region.
\end{proposition}

\begin{protocol}[AI Auditor Operational Workflow]\label{prot:ai_audit}
\begin{enumerate}[label=(\arabic*),itemsep=3pt]
    \item \textbf{Symptom space embedding:} Embed the patient's free-text description into a unified semantic vector space $\mathcal{V} \subset \R^{d_{\text{emb}}}$.
    \item \textbf{Diagnostic trajectory tracking:} For each consultation, record the complete trajectory: initial description $s_0$ $\to$ physician follow-up question $q_1$ $\to$ supplementary description $s_1$ $\to$ ... $\to$ diagnosis $d$.
    \item \textbf{Gradient computation:} In the semantic space, compute the change in diagnostic direction at each step:
    \begin{equation}
        \psi_{\text{grad}} = \frac{1}{|\mathcal{P}|} \sum_{i=1}^{|\mathcal{P}|-1} \|\nabla_s F(s_i) - \nabla_s F(s_{i-1})\|
    \end{equation}
    where $\mathcal{P}$ is the diagnostic trajectory.
    \item \textbf{Misalignment score:} $\psi = \psi_{\text{grad}} \cdot \psi_{\text{consistency}} \cdot \psi_{\text{outcome}}$, where:
    \begin{itemize}
        \item $\psi_{\text{grad}}$: gradient smoothness of the diagnostic trajectory
        \item $\psi_{\text{consistency}}$: stability of diagnosis across multiple consultations for the same patient
        \item $\psi_{\text{outcome}}$: whether treatment outcomes are consistent with the diagnosis (if available)
    \end{itemize}
    \item \textbf{Alert:} When $\psi > \tau_{\text{warn}}$, trigger a gauge misalignment alert, recommending initiation of multi-expert consultation or coordinate system re-declaration.
\end{enumerate}
\end{protocol}

\subsection{AI Audit vs.\ AI Diagnosis---The Critical Distinction}

\begin{center}
\begin{tabular}{p{4cm} p{5cm} p{5cm}}
\toprule
& \textbf{AI Diagnosis} & \textbf{AI Audit} \\
\midrule
Role & Another ``physician''---holds its own coordinate system & Meta-observer---holds no coordinate system \\
Output & Diagnosis $d$ & Misalignment score $\psi$ \\
Error mode & Its own diagnosis can be wrong---and the error is undetectable & Does not diagnose---only detects inconsistency \\
Relationship to existing system & Competitive---replaces physicians & Collaborative---augments physicians \\
Gauge dependence & Has its own implicit gauge---a new misalignment problem & Detects others' gauge misalignment \\
Patient autonomy & Patient is input---AI is decision-maker & Patient is a coordinate system holder---AI is an auditor \\
\bottomrule
\end{tabular}
\end{center}

\honestcrit{Current AI medical diagnosis research mania overlooks this fundamental distinction. Everyone asks ``can AI diagnose more accurately than physicians''---this is the wrong question. The right question is ``can AI detect when the physician and patient are not aligned.'' The former makes AI a stronger coordinate system holder (creating new misalignment problems). The latter makes AI a coordinate-system-neutral auditor (solving misalignment problems).}

% ================================================================
\section{Complete Mathematical Framework of the Diagnostic Process}
\label{sec:framework}
% ================================================================

\subsection{Differential Geometric Formulation of Diagnosis}

We reformulate the diagnostic operation as a parallel transport problem in differential geometry:

\begin{definition}[Diagnosis as Parallel Transport]\label{def:parallel_transport}
Let $\mathcal{M}_{\pt}$ be the patient symptom space (a Riemannian manifold with metric $g_{\pt}$), and $\mathcal{M}_{\med}$ be the physician knowledge space (a Riemannian manifold with metric $g_{\med}$). The diagnostic operation is: choose a connection $\nabla$, construct a gauge transformation (parallel transport) $P_{\gamma}$ from $\mathcal{M}_{\pt}$ to $\mathcal{M}_{\med}$, and move the patient's state vector $v_{\pt}$ into the physician's space along the diagnostic trajectory $\gamma$.

\textbf{Misdiagnosis} corresponds to: the chosen connection $\nabla$ is inconsistent with the ``natural'' correspondence between the two manifolds---i.e., the result of parallel transport depends on the path choice (nonzero curvature) or parallel transport is globally inconsistent (nontrivial holonomy group).
\end{definition}

\begin{theorem}[Diagnostic Holonomy Theorem]\label{thm:holonomy}
Let physicians $D_1$ and $D_2$ start from the same patient description $s_0$ and follow different diagnostic trajectories (asking different questions, ordering different tests) to reach their respective diagnoses $d_1$ and $d_2$. If the curvature tensor $R \neq 0$ of the gauge connection $\nabla$ of doctor-patient communication, then:
\begin{equation}
    d_1 \neq d_2 \quad \text{even in the limit } M \to \infty \text{ of perfect knowledge}
\end{equation}
This means \textbf{diagnostic disagreement can be a holonomy effect caused by curvature---not anyone's fault}.
\end{theorem}

\begin{proof}[Proof Sketch]\rigorSketch
In Riemannian geometry, the holonomy group $\text{Hol}(\nabla)$ describes the ``deficit'' of parallel transport along closed loops---the difference between a vector after going around a loop and the original vector. Different questioning trajectories in the diagnostic process constitute different ``paths,'' and the final diagnosis is the cumulative gauge transformation along the path. If $\text{Hol}(\nabla)$ is nontrivial (curvature nonzero), then the endpoints of two paths differ. This requires no physician to ``make a mistake''---it is a geometric necessity.
\end{proof}

\honestcrit{This theorem has a philosophical consequence: it shows that the assumption ``there exists a unique correct diagnosis'' is equivalent to assuming that the gauge curvature of doctor-patient communication is zero---i.e., a globally trivial parallel transport exists between $\mathcal{M}_{\pt}$ and $\mathcal{M}_{\med}$. This is an extremely strong assumption that almost certainly does not hold in reality. In reality, the correspondence between the two manifolds is local, approximate, and dependent on specific symptom combinations---this is curvature being nonzero.}

\subsection{Variational Formulation of Gauge-Fixing}

\begin{theorem}[Optimal Gauge-Fixing as a Variational Problem]\label{thm:variational_gauge}
The optimal gauge-fixing $\mathbf{g}^*$ is the solution to the following variational problem:
\begin{equation}
    \mathbf{g}^* = \argmin_{\mathbf{g}} \mathbb{E}_{s \sim \mathcal{D}_{\pt}} \left[ \mathcal{L}_{\text{align}}(T_{\mathbf{g}}(s), s_{\text{med-ref}}) \right] + \lambda \cdot \|\mathbf{g}\|_{\text{reg}}
\end{equation}
where:
\begin{itemize}
    \item $\mathcal{L}_{\text{align}}$ is the gauge alignment loss---measuring the distance between the transformed patient description $s$ under $T_{\mathbf{g}}$ and a medical reference standard $s_{\text{med-ref}}$
    \item $s_{\text{med-ref}}$ is a ``zero-gauge'' reference standard constructed through multi-expert consensus (the $M \to \infty$ limit)
    \item $\lambda \cdot \|\mathbf{g}\|_{\text{reg}}$ is a regularization term---penalizing gauges that deviate excessively from the ``null hypothesis'' ($\mathbf{g}=0$ corresponds to ``I faithfully translate the patient's description'')
\end{itemize}
\end{theorem}

\begin{remark}
The choice of $\lambda$ is a political-ethical choice, not a technical choice. $\lambda \to \infty$ means the physician makes no adjustment to the patient's description ($\mathbf{g}=0$)---the extreme of ``what the patient says is true.'' $\lambda \to 0$ means the physician can make arbitrary adjustments to the patient's description---the extreme of ``doctor knows best.'' The optimal $\lambda$ must be determined by social consensus---it is not solvable by an algorithm.
\end{remark}

% ================================================================
\section{Potential Surface Analysis of Special Diagnostic Scenarios}
\label{sec:special_scenarios}
% ================================================================

\subsection{Emergency Medicine: The Most Extreme Case of Potential Surface Misalignment}

Doctor-patient interaction in the emergency department is the extreme aggregation of all potential surface misalignment problems:

\begin{enumerate}[label=(\arabic*),itemsep=2pt]
    \item \textbf{Time compression:} Time for coordinate system declaration is zero---the physician must complete gauge-fixing within seconds
    \item \textbf{Information loss:} The patient may be unconscious or unable to communicate---$\mathcal{S}_{\pt}$ is completely unobservable
    \item \textbf{High-stakes asymmetry:} The cost of missed diagnosis far exceeds the cost of over-diagnosis---$\G^{\text{risk}}$ is pushed to the extreme $r \to 1$
    \item \textbf{High-noise environment:} Pain, fear, and confusion drive the signal-to-noise ratio of patient description extremely low
\end{enumerate}

Under the SCX framework, the optimal diagnostic strategy in the emergency department is not ``get the best physician''---it is \textbf{have the most diverse observers conduct independent assessments in the shortest time}:

\begin{corollary}[Optimal Multi-Observer Strategy for Emergency Medicine]\label{cor:emergency}
Within a time window $\Delta t$, the accuracy-maximizing strategy for emergency diagnosis is:
\begin{equation}
    \max_{\{M, \Delta t_m\}} \Pbb(\text{correct}) \quad \text{s.t.} \quad \sum_{m=1}^{M} \Delta t_m \leq \Delta t
\end{equation}
The optimal solution is not $M=1$ (one physician using all the time), but approximately $M \approx \sqrt{\Delta t / \tau_0}$ observers each spending $\tau_0$ (the minimum effective observation time) conducting independent assessments and then voting.

\textbf{Translation:} Three physicians each looking for 5 minutes $>$ one physician looking for 15 minutes. Because the benefit of independent gauge sampling ($e^{-2M\Delta^2}$) exceeds the benefit of a single physician spending more time when $M$ is small.
\end{corollary}

\subsection{Psychiatry: $\mathcal{S}_{\pt}$ Is Itself the Diagnostic Object}

Psychiatry is the only medical domain where $\mathcal{S}_{\pt}$ is not a signal to be translated but the diagnostic object itself. In psychiatry:

\begin{itemize}
    \item The patient's potential surface $\mathcal{S}_{\pt}$ \textbf{is} the ``organ'' to be evaluated
    \item The physician's coordinate system $\mathcal{S}_{\med}$ is a \emph{description} of $\mathcal{S}_{\pt}$---not an ``objective'' reference independent of $\mathcal{S}_{\pt}$
    \item Hence, the risk of gauge misalignment is highest---because the two coordinate systems are \textbf{conceptually} inequivalent
\end{itemize}

\begin{proposition}[The Psychiatric Gauge Paradox]\label{prop:psychiatry_paradox}
In psychiatric diagnosis, the deviation between $\mathcal{S}_{\pt}$ and $\mathcal{S}_{\med}$ is ineliminable---not a technical problem, but a conceptual problem. $\mathcal{S}_{\med}$ attempts to classify $\mathcal{S}_{\pt}$ as ``normal'' or ``abnormal,'' but the very definition of ``normal'' and ``abnormal'' is itself a dimension of $\mathbf{g}$---the physician is using their own cultural norms to define normalcy.

Hence the ``ground truth'' of psychiatric diagnosis is, at the deepest level, a \textbf{gauge choice}---not an empirically verifiable fact.
\end{proposition}

\honestcrit{This is not anti-psychiatry---it is a precise mathematical formulation of psychiatric uncertainty. The normativity of psychiatric diagnosis (``this is a disease'' / ``this is not a disease'') is not a medical fact---it is a socially determined gauge-fixing. Acknowledging this does not deny psychiatry's value---it makes psychiatry more honest.}

\subsection{Chronic Pain: Prolonged Conflict Between Two Coordinate Systems}

Chronic pain is the most paradigmatic case of doctor-patient potential surface misalignment:

\begin{itemize}
    \item In the patient's $\mathcal{S}_{\pt}$, the pain is \textbf{real, persistent, and has a definite location}.
    \item In the physician's $\mathcal{S}_{\med}$, if no organic lesion can be found---imaging normal, laboratory normal, physical examination normal---then the pain ``should not exist.''
    \item The conflict arises because $\mathcal{S}_{\med}$ lacks a coordinate axis that can accommodate ``pain without organic lesion''---this is not the physician's fault, but a limitation of the current medical paradigm.
\end{itemize}

\begin{definition}[Potential Surface Conflict in Chronic Pain]\label{def:chronic_pain}
Let the patient's pain state be $\omega_{\pt}^{\text{pain}} \in \Omega_{\pt}$---a point of high potential on $\mathcal{S}_{\pt}$. The physician searches for a corresponding point $\omega_{\med}^{\text{pain}}$ in $\Omega_{\med}$, but this point has no correspondence in current medical knowledge---$\mathcal{S}_{\med}(\omega_{\med}^{\text{pain}})$ is undefined. The physician faces two choices:
\begin{enumerate}[label=(\roman*)]
    \item Admit that $\omega_{\med}^{\text{pain}}$ does not exist---``I don't know why you're in pain'' (honest but no treatment)
    \item Map $\omega_{\pt}^{\text{pain}}$ to the nearest defined point in $\Omega_{\med}$---``you might have mild arthritis'' / ``it could be fibromyalgia'' / ``it could be psychological'' (has a diagnosis but may be wrong)
\end{enumerate}
Both choices are consequences of gauge misalignment---the root problem is that there exists a class of states in $\Omega_{\pt}$ with no correspondence in $\Omega_{\med}$.
\end{definition}

\begin{protocol}[Gauge Alignment Operations for Chronic Pain]\label{prot:chronic_pain}
\begin{enumerate}[label=(\arabic*),itemsep=3pt]
    \item \textbf{Acknowledge misalignment:} The physician explicitly tells the patient: ``Your pain is real---it is real in your coordinate system. My coordinate system currently has no entry that perfectly explains it. This does not mean your pain does not exist.''
    \item \textbf{Dual coordinate system parallelism:} Do not cancel either coordinate system. In $\mathcal{S}_{\med}$, continue searching for treatable organic pathology (if any). In $\mathcal{S}_{\pt}$, use the patient's self-report as the measure of treatment effectiveness.
    \item \textbf{Validate in output space:} Do not pursue diagnostic agreement (because the coordinate systems differ), pursue outcome agreement---whether the patient feels improvement after treatment (reduction of potential on $\mathcal{S}_{\pt}$).
\end{enumerate}
\end{protocol}

% ================================================================
\section{The Equality Principle---$\sum \mathbf{g}_m = \mathbf{0}$ in Healthcare Systems}
\label{sec:equality_medicine}
% ================================================================

\subsection{Potential Surface Analysis of Medical Hierarchy}

\begin{definition}[Medical Potential Hierarchy]\label{def:medical_hierarchy}
The hierarchical system in healthcare constitutes a potential staircase:
\begin{equation}
    \mathcal{S}_{\text{hierarchy}} = [
    \text{Patient} \ll \text{Nurse} \ll \text{Resident} \ll \text{Attending} \ll \text{Chief} \ll \text{Academic Authority}
    ]
\end{equation}
The potential difference $\Delta_{ij}$ between each level defines the difficulty of transmitting information upward and the ease of transmitting orders downward. The patient sits at the bottom of the potential well---information can flow downward to the patient (medical orders), but information flowing upward from the patient to the physician (symptom descriptions) decays severely in the potential gradient.
\end{definition}

\begin{proposition}[Directional Bias of Information Flow]\label{prop:info_asymmetry}
In the medical potential hierarchy, information fidelity is directional:
\begin{equation}
    \text{Physician} \to \text{Patient (orders)}: \text{fidelity} \approx 1
\end{equation}
\begin{equation}
    \text{Patient} \to \text{Physician (symptom description)}: \text{fidelity} = 1 - \sum_{i} \Delta_i \cdot \eta_i
\end{equation}
where $\Delta_i$ is the potential difference between the patient and the $i$-th level physician, and $\eta_i$ is the information decay coefficient at each level. The more intermediate levels, the more severe the decay of the patient's voice.
\end{proposition}

\honestcrit{The hierarchical system of current healthcare is not the optimal structure for knowledge transmission---it is the historical inertia of power distribution. A healthcare system with fully smooth potential surfaces (universal equality of speaking weight) is not necessary---but a system in which the patient's voice is systematically decayed mathematically produces more potential surface misalignment. This is not politics---it is information theory.}

\subsection{Shared Decision Making---Institutional Implementation of Gauge-Fixing}

\begin{definition}[Shared Decision Making as a Gauge-Fixing Mechanism]\label{def:shared_decision}
Shared Decision Making (SDM), under the SCX framework, is mathematically defined as:
\begin{equation}
    \mathbf{g}_{\text{fixed}} = \frac{\mathbf{g}_{\med} + \mathbf{g}_{\pt}}{2} \quad \text{(equal-weight average)}
\end{equation}
or more generally:
\begin{equation}
    \mathbf{g}_{\text{fixed}} = w_{\med} \cdot \mathbf{g}_{\med} + w_{\pt} \cdot \mathbf{g}_{\pt}, \quad w_{\med} + w_{\pt} = 1
\end{equation}
where $w_{\pt} > 0$---i.e., the patient's coordinate system \textbf{must} have nonzero weight.
\end{definition}

\begin{theorem}[Mathematical Guarantee That SDM Reduces Misdiagnosis]\label{thm:sdm}
When gauge-fixing uses equal-weight averaging ($w_{\pt} = w_{\med} = \frac{1}{2}$), the upper bound on worst-case diagnostic distance is strictly less than the upper bound under physician-unilateral gauge-fixing ($w_{\pt} = 0$):
\begin{equation}
    \max_{s \in \Sigma} \|F_{\mathbf{g}_{\text{fixed}}^{\text{SDM}}}(s) - F_{\mathbf{g}^{*}}(s)\| < \max_{s \in \Sigma} \|F_{\mathbf{g}_{\text{fixed}}^{\text{doc-only}}}(s) - F_{\mathbf{g}^{*}}(s)\|
\end{equation}
unless $\mathcal{S}_{\pt}$ and $\mathcal{S}_{\med}$ completely coincide (in which case they are equivalent). In all nontrivial cases, SDM strictly reduces worst-case error.
\end{theorem}

\begin{proof}[Proof Sketch]\rigorSketch
By the triangle inequality: the worst-case deviation of the average point is strictly less than the worst-case deviation of either endpoint, unless the two coordinate systems yield the same diagnostic function (in which case the deviations are equal). In medical reality, $\mathcal{S}_{\pt} \neq \mathcal{S}_{\med}$ almost always holds (nontrivial case), so the strict inequality holds.
\end{proof}

\subsection{Nine Medical Equality Principles}

The following nine items constitute all operational constraints of the Equality Principle in healthcare systems---each is a direct corollary of the preceding theorems:

\begin{enumerate}[label=\textbf{(\arabic*)}]
    \item \textbf{Medical hierarchy necessarily decays the patient's voice.} The more intermediate levels, the longer the information channel from patient to decision-making physician, and the greater the decay. This is not ``bad attitude''---it is channel physics. Flattening the diagnostic structure (reducing intermediate levels, having patients directly face decision-making physicians) is a mathematical necessity for reducing information decay.

    \item \textbf{Patient autonomy is not ``granted''---it is ``acknowledged.''} $w_{\pt} > 0$ is not the physician generously ceding power---it is $\sum \mathbf{g}_m = \mathbf{0}$ demanding that the patient coordinate system must participate in gauge-fixing. A physician ``not respecting patient autonomy'' is mathematically equivalent to ``claiming one's own coordinate system is the only valid one''---this claim is mathematically invalid.

    \item \textbf{Gauge diversity in multi-expert consultation is a core value.} The goal of consultation is not to get all experts to ``unify their thinking''---it is to preserve their independent gauge diversity. If consultation becomes ``unify to the chief's gauge,'' it loses the exponential benefit of $e^{-2M\Delta^2}$. $M$ physicians with identical gauges = 1 physician.

    \item \textbf{The optimal role of AI is as an auditor---not as a physician replacement.} Making AI another ``physician'' (holding its own coordinate system) creates a new potential surface misalignment problem. Making AI a potential surface auditor (detecting misalignment between others' coordinate systems) solves the potential surface misalignment problem. The right question is ``is this diagnostic trajectory consistent between the patient coordinate system and the doctor coordinate system''---not ``is the AI's diagnosis correct.''

    \item \textbf{Attitude auditing must become a routine indicator of healthcare quality.} A physician's implicit gauge bias---gender, age, racial, class bias---is not a marker of a ``bad person,'' but the degree of freedom of an unfixed gauge. But a degree of freedom is not harmless simply because it is ``unconscious.'' Attitude auditing (Prop.~\ref{prop:attitude}) provides a non-accusatory, statistically data-driven method of bias detection.

    \item \textbf{Patients with chronic pain / functional disorders should not be questioned about the ``reality'' of their experience.} When a state exists in the patient's $\mathcal{S}_{\pt}$ with no correspondence in the physician's $\mathcal{S}_{\med}$, the question is not ``is the patient's description truthful''---it is ``does a correspondence exist between the two coordinate systems.'' The latter is a medical knowledge boundary problem---not a patient credibility problem.

    \item \textbf{Accountability for diagnostic error should distinguish knowledge error from gauge error.} Knowledge error (not knowing this disease) concerns room for improvement. Gauge error (knowing this disease but failing to correctly map the patient's description in one's own coordinate system) concerns coordinate system misalignment. Confusing these two types of error leads to physicians being blamed for gauge errors (``how could you miss such an obvious diagnosis'')---when in fact the blame should lie with a process lacking a gauge-fixing mechanism.

    \item \textbf{The ``normal''/``abnormal'' boundary in psychiatric diagnosis requires explicit declaration of its gauge dependence.} Psychiatric diagnostic criteria (e.g., DSM, ICD) define a series of ``disorders.'' From the perspective of potential surface geometry, every diagnostic criterion is a gauge choice---it draws a line on the patient's $\mathcal{S}_{\pt}$, with ``normal'' on one side and ``disorder'' on the other. The position of this line is not naturally given---it is the result of social negotiation. Psychiatry should declare its gauge dependence more honestly than any other medical domain.

    \item \textbf{$\sum \mathbf{g}_m = \mathbf{0}$ is the only mathematical formulation of healthcare equity.} The mathematically equivalent form of health equity is: no person's coordinate system is systematically excluded from the gauge-fixing process. If the $\mathcal{S}_{\pt}$ of a certain group (defined by race, gender, class, language) is never incorporated into the gauge-fixing of diagnostic norms, then that group faces $\mathbb{E}[\Misdiagnose] > \mathbb{E}[\Misdiagnose]_{\text{population}}$---this is systematic misalignment---this is inequity.
\end{enumerate}

\textbf{One sentence:} Healthcare equity is not ``everyone receives the same treatment''---it is ``everyone's coordinate system participates in diagnostic gauge-fixing with equal weight.'' $\sum \mathbf{g}_m = \mathbf{0}$.

% ================================================================
\section{Honest Strike---Limitations and Reflections on This Paper Itself}
\label{sec:honest_strike}
% ================================================================

\honestcrit{The following content is a Honest Strike against this paper itself. Any theory claiming to have discovered a ``fundamental problem'' must first endure the same scrutiny.}

\subsection{Strike 1: The Cost of Formalization}

Medical diagnosis is a process that involves the ineffable---the patient's experience of suffering, the physician's clinical intuition, the subtle nonverbal interaction between the two. Formalizing these as parallel transport on manifolds and variational principles on gauge groups may capture the skeleton of the structure but lose the flesh of the content.

\textbf{Response:} Formalization is not intended to replace clinical judgment---it is intended to reveal the hidden mathematical structure within clinical judgment. Just as gauge theory in physics is not intended to replace physical intuition---it is intended to reveal why physical intuition is sometimes right and sometimes wrong. The value of this paper's mathematical framework lies not in its being ``more precise''---it lies in explaining \textbf{why} even the most excellent clinicians make systematic errors.

\subsection{Strike 2: Patients Cannot Always Declare Their Coordinate System}

This paper assumes the patient can declare their coordinate system. But many patients---due to disease state (unconsciousness, cognitive impairment), age (children, very elderly), language barriers, or educational limitations---cannot complete a coordinate system declaration. In these cases, the framework of this paper cannot be directly applied.

\textbf{Response:} This is a genuine limitation. In these cases, gauge-fixing must be completed by proxy---family members, guardians, patterns from prior medical records. Proxy gauge introduces a new dimension of misalignment: the proxy's $\mathbf{g}_{\text{proxy}}$ may be inconsistent with the patient's $\mathcal{S}_{\pt}$. This paper acknowledges this additional error source in the chain and flags it as an Open Problem.

\subsection{Strike 3: The Resource Cost of Multi-Expert Consultation}

Theorem~1 says the consensus of $M$ physicians converges at exponential rate. But the time and cost of $M$ physicians is also $O(M)$. In resource-limited healthcare systems, fully independent consultation with $M=10$ is infeasible.

\textbf{Response:} This paper does not advocate that every patient needs $M=10$ consultation. It advocates that: \textbf{when diagnostic uncertainty is high} ($\Delta$ is small---the disease signal is weak), increasing $M$ is the optimal resource allocation strategy. In most routine diagnoses, $\Delta$ is large enough that $M=1$ or $M=2$ (second opinion) suffices. But in difficult cases, rare diseases, or when diagnostic disagreement has already emerged---Theorem~1 tells you: adding one independent physician is more effective than letting the existing physician think for 10 more minutes.

\subsection{Strike 4: Training Data Bias in the AI Auditor}

The AI potential surface auditor is itself trained on data---data that may already contain systematic gauge bias (e.g., historical records of certain populations being systematically misdiagnosed). If the auditor learns these biases, its ``misalignment score'' may not detect misalignment---it may reinforce misalignment.

\textbf{Response:} This is the most important Honest Strike in the design of the AI auditor. The solution involves two aspects: (1) debiasing training data---using multi-source data, balancing population representation; (2) transparency of the auditor itself---the auditor's internal decision pathways must be interpretable, so that its own gauge can be independently audited. Although the AI auditor does not hold a diagnostic coordinate system, it does hold a definition of ``what constitutes misalignment''---this definition itself is also a gauge choice. This paper acknowledges this meta-level misalignment risk and flags it as a continuous calibration requirement for the AI auditor.

\subsection{Strike 5: The Non-Reformability of Current Healthcare Systems}

The set of operations proposed in this paper---coordinate system declaration, independent multi-expert consultation, attitude auditing, AI auditing---may have implementation costs so high as to be unrealistic in current healthcare systems. A 15-minute outpatient slot is insufficient even for basic history-taking, let alone for coordinate system declaration.

\textbf{Response:} This paper is theoretical work, not policy recommendation. It does what theory does---reveal deep structure, explain why the current healthcare system mathematically \emph{must} produce misdiagnosis. How to change the system is a political/economic/organizational problem, not a mathematical problem. But the first step of change is understanding---understanding that misdiagnosis is not individual failure, but systemic inevitability. If this paper can make one policymaker, one hospital administrator, one medical educator pause before a decision and say---``wait, this process has no gauge-fixing''---it has achieved its purpose.

\subsection{Strike 6: The Validity Boundary of the Patient Coordinate System}

This paper assumes the patient's coordinate system $\mathcal{S}_{\pt}$ is ``valid.'' But some patients do lie (insurance fraud, drug-seeking), do exaggerate, do misjudge their own bodies. Fully respecting the patient coordinate system ($w_{\pt} = 1$) means accepting all self-reports as true---which is dangerous in some situations.

\textbf{Response:} This paper does not advocate $w_{\pt} = 1$ (the patient is absolutely correct). It advocates $w_{\pt} > 0$ (the patient coordinate system has nonzero weight). Determining the optimal value of $w_{\pt}$ is an empirical question---it depends on the reliability of patient self-report in specific clinical scenarios. In scenarios with objective verification means (e.g., radiographically confirmable fracture), $w_{\pt}$ can be very low. In scenarios without objective verification means (e.g., pain, fatigue, mood), $w_{\pt}$ must be higher---not because the patient is ``more credible,'' but because \textbf{no other coordinate system is available}. Rejecting the only available reference frame when no other exists---however imperfect it may be---is not prudence, it is ineffectiveness.

\subsection{Strike 7: Empirical Testability of the Holonomy Theorem}

The Diagnostic Holonomy Theorem (Theorem~\ref{thm:holonomy}) asserts that nonzero curvature leads to path-dependent diagnostic disagreement. But how can this curvature be measured in practice? The diagnostic process is unlike a physics experiment---you cannot ``put the same patient back'' and let another physician re-interview them. We can only ever see the result of one diagnostic trajectory.

\textbf{Response:} This is an Open Problem. Possible approaches include:
\begin{itemize}
    \item \textbf{Simulated patients:} Train standardized patients to present the same symptom set, and let different physicians independently interview them---this is the closest to a ``controlled experiment.''
    \item \textbf{Chart review:} From large volumes of medical records, match consultation pairs with ``similar initial descriptions but different questioning trajectories,'' and compare the consistency of final diagnoses---this is an observational approximation.
    \item \textbf{AI simulation:} Use language models to simulate different physicians' questioning trajectories when faced with the same patient description---this is a synthetic data approach, useful for generating hypotheses but not for confirmation.
\end{itemize}

This paper flags these methods as future work, while acknowledging that the Diagnostic Holonomy Theorem is, at this stage, more of a conceptual framework---it points to a phenomenon that should exist (path-dependent diagnostic disagreement), but rigorous empirical verification requires methods that transcend current medical data infrastructure.

% ================================================================
\section{Potential Surface Restatement of Related Work}
\label{sec:related_work}
% ================================================================

The following restates classic concepts from the medical domain in SCX potential surface language:

\begin{center}
\begin{tabular}{p{3.5cm} p{4cm} p{6cm}}
\toprule
\textbf{Traditional Concept} & \textbf{SCX Restatement} & \textbf{Mathematical Formulation} \\
\midrule
Second opinion & Gauge sampling by a second observer & $M \to 2$ special case of Theorem~1 \\
\addlinespace
Shared decision making & Joint optimization of gauge-fixing & $\mathbf{g}_{\text{fixed}} = w_{\med}\mathbf{g}_{\med} + w_{\pt}\mathbf{g}_{\pt}$ \\
\addlinespace
Evidence-based medicine & Anchoring $\mathcal{S}_{\med}$ to population data & $\mathcal{S}_{\med} \approx \mathbb{E}_{\text{pop}}[\mathcal{S}]$ \\
\addlinespace
Patient-reported outcome measures (PROMs) & Direct measurement in $\mathcal{S}_{\pt}$ & $\Delta\mathcal{S}_{\pt}(t)$ as efficacy metric \\
\addlinespace
Clinical practice guidelines & Population approximation of gauge-fixing & $\mathbf{g}_{\text{fixed}}^{\text{pop}}$ as default gauge \\
\addlinespace
Differential diagnosis & Searching multiple minima on $\mathcal{S}_{\med}$ & $\{d: \mathcal{S}_{\med}(\omega_d) < \tau\}$ \\
\addlinespace
Overdiagnosis & Risk gauge pushed to $r \to 1$ & $\G^{\text{risk}}$ unfixed $\to$ diagnostic inflation \\
\addlinespace
Missed diagnosis & Risk gauge pushed to $r \to 0$ & $\G^{\text{risk}}$ unfixed $\to$ diagnostic narrowing \\
\addlinespace
Placebo effect & Therapeutic effect of gauge alignment itself & $\Delta\mathcal{S}_{\pt} > 0$ even if $\Delta\mathcal{S}_{\med} = 0$ \\
\addlinespace
Doctor-patient trust & Shrinking of gauge holonomy group & $\text{Hol}(\nabla) \to 0$ with increasing interaction count \\
\addlinespace
Rare disease diagnostic delay & $\omega_d$ has excessively high potential in $\mathcal{S}_{\med}$ & $\mathcal{S}_{\med}(\omega_d) \gg \tau$, large $M$ needed to exceed noise \\
\bottomrule
\end{tabular}
\end{center}

\subsection{Potential Surface Limitations of Evidence-Based Medicine}

Evidence-Based Medicine (EBM) anchors $\mathcal{S}_{\med}$ to population data: $\mathcal{S}_{\med} \approx \mathbb{E}_{\text{pop}}[\mathcal{S}]$. This is an enormous advance---it liberates physicians from the hallucination that ``personal experience is truth.'' But from the perspective of potential surface geometry, EBM has an under-discussed limitation:

\begin{proposition}[The Implicit Gauge of Evidence-Based Medicine]\label{prop:ebm_gauge}
The ``evidence'' of evidence-based medicine itself carries an implicit gauge---the inclusion/exclusion criteria of clinical trials define a ``standard patient,'' but real patients rarely perfectly match this standard. When applying clinical trial results to a specific patient, the physician \textbf{must perform gauge adjustment}---the magnitude and direction of the adjustment depend on the physician's judgment. EBM claims to have eliminated physician subjectivity---in reality, it merely shifted subjectivity from one dimension (``I think this drug works'') to another dimension (``I think this patient is this far from the clinical trial inclusion criteria'').
\end{proposition}

\honestcrit{This is not a denial of EBM---EBM is one of the most significant advances in the history of medicine. This is pointing out that EBM completed one gauge-fixing (the ``zero point'' of $\mathcal{S}_{\med}$ shifted from individual experience to population statistics), but did not resolve the gauge misalignment between $\mathcal{S}_{\pt}$ and $\mathcal{S}_{\med}$. A good evidence-based physician must ask not only ``what treatment is most effective for the 'average patient'''---but also ``how far is this patient from the 'average patient,' and how should I adjust the coordinates.''}

% ================================================================
\section{Conclusion}
\label{sec:conclusion}
% ================================================================

We have proved a fundamental but unformalized phenomenon in medical diagnosis: physicians and patients live in different coordinate systems, and the diagnostic operation has no predefined gauge-fixing between the two coordinate systems. Misdiagnosis---or more precisely, ``diagnostic uncertainty''---has an ineliminable mathematical boundary within the single-physician, single-coordinate-system framework, given by the Honest Person Theorem (Theorem~3).

Multi-expert consultation ($M > 1$) provides a gauge-fixing mechanism---the consensus detection of $M$ independently gauged physicians reduces misdiagnosis at the exponential convergence rate $e^{-2M\Delta^2}$ (Theorem~1). This is not the intuition that ``many hands make light work''---it is a mathematical guarantee.

Patient autonomy acquires a new mathematical foundation in this paper's framework: it is not an ethical appeal---it is the inevitable consequence of $\sum \mathbf{g}_m = \mathbf{0}$. The patient's coordinate system $\mathcal{S}_{\pt}$ must participate in the diagnostic gauge-fixing process with nonzero weight---not as a gesture of ``respect,'' but as a mathematical operation for reducing information loss.

The correct role of AI-assisted diagnosis is not to become another ``physician'' (holding its own coordinate system---creating new misalignment problems), but to become a \textbf{potential surface auditor}---detecting misalignment between physician and patient coordinate systems without holding its own coordinate system.

We have provided three operational protocols (coordinate system declaration, attitude auditing, AI potential surface auditing), nine Medical Equality Principle articles, and a Honest Strike section (scrutinizing the limitations of this paper itself). We have developed a differential geometric formulation of diagnosis, defined the gauge group $\G$ (comprising translation, scaling, prior, and risk subgroups), and established a variational principle for optimal gauge-fixing.

\textbf{Final message:} Align the coordinate systems before diagnosing. This is not just good medical practice---it is a mathematical necessity. If a physician does not acknowledge that their own coordinate system is merely one coordinate system---not the only coordinate system---then their diagnosis necessarily contains an error component that no amount of training can eliminate. This component does not reflect their incompetence---it reflects an incomplete operation: $\mathbf{g}_{\med}$ and $\mathbf{g}_{\pt}$ have not yet been fixed to the same reference.

\textbf{Final sentence:} The patient says ``I hurt''---in the patient's coordinate system, this is true. The physician says ``I cannot find a cause''---in the physician's coordinate system, this is also true. The conflict between two truths does not require one party to be false. It requires not a smarter physician---it requires a gauge-fixing that aligns the two coordinate systems. $\sum\mathbf{g}_m = \mathbf{0}$.

% ================================================================
% Acknowledgments
% ================================================================
\vspace{10pt}
\noindent\textbf{Acknowledgments:} This work was completed independently under the SCX framework. Gratitude to all honest physicians---you strive under impossible conditions to align two incommensurable coordinate systems; the theorems are not a critique of you. Gratitude to all patients---your coordinate systems are real, even when medicine cannot yet read them.

\vspace{10pt}
\noindent\textbf{Author's Statement:} The author publishes under the name SCX. This work is released on GitHub for immediate open access. No institutional affiliation is claimed. The theorems stand on their own. 

This work does not constitute medical advice. The mathematical framework presented here is descriptive and analytical---it reveals structural properties of the diagnostic process, but clinical decisions must be made by qualified medical professionals using their full judgment.

% ================================================================
% References
% ================================================================
\begin{thebibliography}{99}

\bibitem{scx_moe_gauge}
SCX,
``Potential Surface Misalignment---Gauge Freedom in Multi-Expert Routing and MILP Gauge-Fixing,''
GitHub: shuchaoxi/SCX, 2026.

\bibitem{scx_hamiltonian}
SCX,
``Hamiltonian as Audit Condition---Judging Auditability from the Energy Landscape,''
GitHub: shuchaoxi/SCX, 2026.

\bibitem{scx_theory}
SCX,
``Detecting Label Noise by State-Conditioned eXpertise,''
GitHub: shuchaoxi/SCX, 2026.

\bibitem{scx_thm3}
SCX,
``Honest Person Theorem: Single Observer Cannot Distinguish Noise, Bias, Learnable Difficulty, and Honest Error Without Additional Assumptions,''
GitHub: shuchaoxi/SCX, 2026.

\bibitem{scx_yajie}
SCX,
``Yajie Protocol: Honest Nash Equilibrium in Multi-Expert Consensus,''
GitHub: shuchaoxi/SCX, 2026.

\bibitem{scx_galois}
SCX,
``Galois Unsolvability Theorem: The Death Sentence of Recursive Distillation,''
GitHub: shuchaoxi/SCX, 2026.

\bibitem{who_patient_safety}
World Health Organization,
``Global Patient Safety Action Plan 2021--2030,''
WHO, 2021.

\bibitem{balogh2015}
E. Balogh, B. T. Miller, and J. Ball (eds.),
``Improving Diagnosis in Health Care,''
National Academies Press, 2015.

\bibitem{singh2014}
H. Singh, A. N. D. Meyer, and E. J. Thomas,
``The frequency of diagnostic errors in outpatient care,''
\textit{BMJ Quality \& Safety}, 2014.

\bibitem{graber2012}
M. L. Graber, R. M. Wachter, and C. K. Cassel,
``Bringing diagnosis into the quality and safety equations,''
\textit{JAMA}, 2012.

\bibitem{sackett1996}
D. L. Sackett, W. M. C. Rosenberg, J. A. M. Gray, R. B. Haynes, and W. S. Richardson,
``Evidence based medicine: what it is and what it isn't,''
\textit{BMJ}, 1996.

\bibitem{elwyn2012}
G. Elwyn, D. Frosch, R. Thomson, et al.,
``Shared Decision Making: A Model for Clinical Practice,''
\textit{Journal of General Internal Medicine}, 2012.

\bibitem{charles1997}
C. Charles, A. Gafni, and T. Whelan,
``Shared decision-making in the medical encounter,''
\textit{Social Science \& Medicine}, 1997.

\bibitem{topol2019}
E. Topol,
``High-performance medicine: the convergence of human and artificial intelligence,''
\textit{Nature Medicine}, 2019.

\bibitem{rajpurkar2022}
P. Rajpurkar, E. Chen, O. Banerjee, and E. J. Topol,
``AI in health and medicine,''
\textit{Nature Medicine}, 2022.

\bibitem{obermeyer2019}
Z. Obermeyer, B. Powers, C. Vogeli, and S. Mullainathan,
``Dissecting racial bias in an algorithm used to manage the health of populations,''
\textit{Science}, 2019.

\bibitem{green2021}
A. R. Green, D. R. Carney, D. J. Pallin, et al.,
``Implicit bias among physicians and its prediction of thrombolysis decisions for black and white patients,''
\textit{Journal of General Internal Medicine}, 2007.

\bibitem{hoffman2016}
K. M. Hoffman, S. Trawalter, J. R. Axt, and M. N. Oliver,
``Racial bias in pain assessment and treatment recommendations,''
\textit{PNAS}, 2016.

\bibitem{fitzgerald2017}
C. Fitzgerald and S. Hurst,
``Implicit bias in healthcare professionals: a systematic review,''
\textit{BMC Medical Ethics}, 2017.

\bibitem{armstrong2020}
D. Armstrong,
``The social life of diagnosis: reflections on the DSM and medical classification,''
\textit{Sociology of Health \& Illness}, 2020.

\bibitem{engel1977}
G. L. Engel,
``The need for a new medical model: a challenge for biomedicine,''
\textit{Science}, 1977.

\bibitem{kleinman1988}
A. Kleinman,
\textit{The Illness Narratives: Suffering, Healing, and the Human Condition},
Basic Books, 1988.

\bibitem{canguilhem1991}
G. Canguilhem,
\textit{The Normal and the Pathological},
Zone Books, 1991 (orig. 1943).

\bibitem{foucault1973}
M. Foucault,
\textit{The Birth of the Clinic: An Archaeology of Medical Perception},
Vintage Books, 1973.

\bibitem{hoeffding}
W. Hoeffding,
``Probability Inequalities for Sums of Bounded Random Variables,''
\textit{Journal of the American Statistical Association}, 1963.

\bibitem{nakamura2022}
Y. Nakamura, K. Kiyono, and T. Kudo,
``Geometric deep learning for medical diagnosis,''
\textit{Medical Image Analysis}, 2022.

\bibitem{vaswani2017}
A. Vaswani, N. Shazeer, N. Parmar, et al.,
``Attention Is All You Need,''
\textit{NeurIPS}, 2017.

\bibitem{bahadur1960}
R. R. Bahadur and R. R. Rao,
``On deviations of the sample mean,''
\textit{Annals of Mathematical Statistics}, 1960.

\end{thebibliography}

\end{document}
